% !TEX root = z-main.tex

\section{Justificación}

El modelado tridimensional de las ruinas arqueológicas de Kaminaljuyú constituye un aporte pertinente para la preservación y la divulgación del patrimonio cultural guatemalteco, al ofrecer representaciones digitales que facilitan el acceso, la interpretación y el aprendizaje sobre un sitio cuyo paisaje arqueológico ha sido sustancialmente alterado por la urbanización.

\subsection*{Dimensión patrimonial: pérdida por urbanización y acceso al patrimonio}
Durante la época prehispánica, Kaminaljuyú abarcó aproximadamente cinco kilómetros cuadrados y estuvo conformado por más de 230 montículos. Debido al crecimiento urbano de la ciudad de Guatemala, gran parte del sitio fue destruido o quedó cubierto por construcciones modernas; se estima que hoy se conservan de forma visible alrededor de cuarenta montículos, de los cuales nueve se resguardan en el Parque Arqueológico Kaminaljuyú \cite{HemerotecaPL2016}. Esta reducción del paisaje original limita la posibilidad de observación y estudio in situ, particularmente en estructuras emblemáticas como la Acrópolis (C-II-4) o sectores de La Palangana, donde la erosión y el paso del tiempo dificultan comprender sus volúmenes y relaciones espaciales. En este contexto, la reconstrucción digital aporta una ventana interpretativa que, sin sustituir el trabajo arqueológico, posibilita recuperar visualmente formas, jerarquías y disposiciones que hoy no son evidentes \cite{FundacionWiese2017}. Además, se alinea con iniciativas nacionales de digitalización de bienes culturales impulsadas por el Ministerio de Cultura y Deportes \cite{MCD2024}.

\subsection*{Dimensión educativa: aprendizaje significativo y divulgación accesible}
La integración de modelos \gls{3d} en entornos de realidad aumentada favorece procesos de aprendizaje activo, al permitir la exploración situada de estructuras y la relación entre el espacio físico y la información digital. La literatura señala que la RA incrementa la motivación y mejora la comprensión de contenidos complejos al ofrecer visualizaciones interactivas y contextuales, con efectos positivos en la retención del conocimiento \cite{Barroso2022,Sevilla2017,Handley2024}. En el ámbito universitario y museográfico, estas experiencias facilitan la enseñanza de historia y arqueología al presentar reconstrucciones creíbles, comparaciones antes/después y lecturas volumétricas del sitio; al mismo tiempo, habilitan recursos inclusivos para públicos no especializados mediante dispositivos móviles de amplio acceso. En consecuencia, el proyecto aporta materiales didácticos reutilizables para cursos, visitas guiadas y programas de mediación cultural, contribuyendo a la alfabetización patrimonial y tecnológica de estudiantes y visitantes.

\subsection*{Dimensión tecnológica: documentación 3D y RA como estrategia de preservación digital}
Desde el punto de vista técnico, el uso de Blender para modelado y texturizado, junto con la preparación de activos para su integración en \gls{RA}, promueve buenas prácticas de documentación y preservación digital (formatos interoperables, metadatos, trazabilidad de versiones). La estandarización hacia formatos orientados a tiempo real  facilita la portabilidad entre plataformas y la sostenibilidad a largo plazo de los recursos digitales. Estas prácticas se encuentran en sintonía con los lineamientos internacionales sobre preservación del patrimonio digital, que enfatizan la accesibilidad, autenticidad y continuidad de los contenidos culturales en entornos digitales \cite{UNESCO2003Charter}. En conjunto, el proyecto no solo provee visualizaciones para la difusión, sino que genera insumos técnicos que pueden integrarse en repositorios académicos o institucionales para su conservación y reutilización futura.

\subsection*{Síntesis}
La convergencia de estas tres dimensiones justifica la pertinencia del proyecto: (i) responde a la pérdida material del sitio mediante reconstrucciones verificables; (ii) habilita experiencias educativas accesibles y significativas; y (iii) establece un proceso técnico alineado con la preservación digital y la interoperabilidad, fortaleciendo el ecosistema de investigación y divulgación del patrimonio arqueológico guatemalteco.